% !TeX program = xelatex

\documentclass[aspectratio=169]{beamer}

% ============================================================
% 1. Fonts
% ============================================================
\usepackage{fontspec}

\setmainfont{Microsoft JhengHei}
\setsansfont{Microsoft JhengHei}
\setmonofont{JetBrains Mono}

\XeTeXlinebreaklocale "zh"


% ============================================================
% 2. Mathematics, enumerate
% ============================================================
\usepackage{amsmath}
\usepackage{amssymb}
\usepackage{enumitem}

\setlistdepth{6}
\renewlist{itemize}{itemize}{9}
\renewlist{enumerate}{enumerate}{9}
\setlist[itemize,1]{label=\textbullet}
\setlist[itemize,2]{label=$\circ$}
\setlist[itemize,3]{label=$\ast$}
\setlist[itemize,4]{label=$\cdot$}
\setlist[itemize,5]{label=--}
\setlist[itemize,6]{label=--}
\setlist[enumerate,1]{label=\arabic*.}
\setlist[enumerate,2]{label=\arabic*.}
\setlist[enumerate,3]{label=\arabic*.}
\setlist[enumerate,4]{label=\arabic*.}


% ============================================================
% 3. Images / Tables / Boxes
% ============================================================
\usepackage{graphicx}
\usepackage{booktabs}
\usepackage[most]{tcolorbox}


% ============================================================
% 4. Colors
% ============================================================
\definecolor{myRed}{RGB}{145,90,90}
\definecolor{myGreen}{RGB}{85,115,100}
\definecolor{myBlue}{RGB}{70,100,130}
\definecolor{myOrange}{RGB}{155,115,80}
\definecolor{myPurple}{RGB}{115,100,130}

\definecolor{myGray}{RGB}{242,242,238}
\definecolor{myDarkgray}{RGB}{90,90,88}
\definecolor{myBlack}{RGB}{45,45,43}


% ============================================================
% 5. Beamer Colors
% ============================================================
\setbeamercolor{structure}   {fg=myBlue}
\setbeamercolor{title}       {fg=white, bg=myBlue}
\setbeamercolor{frametitle}  {fg=white, bg=myBlue}
\setbeamercolor{block title} {fg=white, bg=myBlue}
\setbeamercolor{block body}  {fg=myBlack, bg=myGray}
\setbeamercolor{alerted text}{fg=myRed}
\setbeamercolor{example text}{fg=myGreen}


% ============================================================
% 6. Theme
% ============================================================
\usetheme{Madrid}
\usecolortheme{default}

\setbeamertemplate{navigation symbols}{}


% ============================================================
% 7. Footer
% ============================================================
\setbeamertemplate{footline}{
    \leavevmode
    \hbox{
        % Left: Author
        \begin{beamercolorbox}[
            wd=.33\paperwidth,
            ht=2.5ex,
            dp=1ex,
            center
        ]{author in head/foot}
            \insertdate
        \end{beamercolorbox}

        % Center: Title
        \begin{beamercolorbox}[
            wd=.33\paperwidth,
            ht=2.5ex,
            dp=1ex,
            center
        ]{title in head/foot}
            \insertshorttitle
        \end{beamercolorbox}

        % Right: Page Number
        \begin{beamercolorbox}[
            wd=.33\paperwidth,
            ht=2.5ex,
            dp=1ex,
            center
        ]{date in head/foot}
            \insertframenumber{} / \inserttotalframenumber
        \end{beamercolorbox}
    }
}


% ============================================================
% 8. Presentation Information
% ============================================================
\title{簡報模板暨使用範例}
\author{Syu}
\date{September 17, 2026}


% ========================= DocStart =========================
\begin{document}

% ============================================================
% Title Page
% ============================================================
\begin{frame}
    \titlepage
\end{frame}


% ============================================================
% Index
% ============================================================
\begin{frame}{Index}
    \tableofcontents
\end{frame}


% ============================================================
% Section
% ============================================================
\section{基本語法}


% ============================================================
% Basic Text
% ============================================================
\begin{frame}{文字}
    一般文字。
    
    \bigskip
    
    \textbf{粗體文字}
    
    \textit{斜體文字}
    
    \alert{警告或重點文字}

    \href{https://github.com/shihyousyu}{超連結}
    
\end{frame}


% ============================================================
% Itemize
% ============================================================
\begin{frame}{項目列表}
    \begin{itemize}
    
        \item 第一層
        
        \begin{itemize}
        
            \item 第二層
            
            \begin{itemize}
            
                \item 第三層
                
            \end{itemize}
            
        \end{itemize}
        
        \item 第二個項目
        
    \end{itemize}
\end{frame}


% ============================================================
% Enumerate
% ============================================================
\begin{frame}{編號列表}
    \begin{enumerate}
    
        \item 第一項
        \item 第二項
        \item 第三項
        
    \end{enumerate}
\end{frame}


% ============================================================
% Blocks
% ============================================================
\begin{frame}{Block}
    \begin{block}{一般 Block}
        這是一個一般的 Block。
    \end{block}
    
    \begin{alertblock}{Alert Block}
        這是一個警告 Block。
    \end{alertblock}

    \begin{exampleblock}{Example Block}
        這是一個範例 Block。
    \end{exampleblock}
\end{frame}


% ============================================================
% Mathematics
% ============================================================
\begin{frame}{數學公式}
    行內公式：
    $E = mc^2$
    \bigskip

    顯示公式：
    \[
        E = mc^2
    \]

    \begin{equation}
        a^2 + b^2 = c^2
    \end{equation}
\end{frame}


% ============================================================
% Columns
% ============================================================
\begin{frame}{雙欄}
    \begin{columns}

        % ========================
        % Left
        % ========================
        \begin{column}{0.48\textwidth}
            \begin{block}{左側}
                左側內容。
            \end{block}
        \end{column}


        % ========================
        % Right
        % ========================
        \begin{column}{0.48\textwidth}
            \begin{block}{右側}
                右側內容。
            \end{block}
        \end{column}
    \end{columns}
\end{frame}


% ============================================================
% Table
% ============================================================
\begin{frame}{表格}
    \begin{table}
        \centering

        \begin{tabular}{lcc}

            \toprule

            Language & Level & Usage \\

            \midrule

            C/C++ & High & System \\

            Python & High & Data \\

            Java & Medium & Backend \\

            \bottomrule
            
        \end{tabular}
    \end{table}
\end{frame}


% ============================================================
% Image
% ============================================================
\begin{frame}{圖片}

    \centering
    
    \begin{figure}
        \centering
        \includegraphics[width=.5\linewidth]{you don't have to be great to start, but you have to start to be great.png}
        \caption{Example}
        \label{fig:placeholder}
    \end{figure}

\end{frame}


% ============================================================
% Tcolorbox
% ============================================================
\begin{frame}{自訂 Box}
    \begin{tcolorbox}[
        colback=myGray,
        colframe=myBlue,
        title=Custom Box
    ]
        這是一個使用 tcolorbox 製作的自訂區塊。
    \end{tcolorbox}
\end{frame}


% ============================================================
% Final
% ============================================================
\begin{frame}{結束}

    \centering

    \Huge
    Thank You

    \bigskip

    \large
    Questions?

\end{frame}


\end{document}